\chapter{Westernization \& Sanskritization}

\section{Westernization: Three Levels of Change}

Westernization (M. N. Srinivas's idea) is a process of change at \textbf{three levels} over 150 years of British rule --- remember \textbf{IIT}: \textbf{Institutional, Ideological, Technological}.
\begin{itemize}
\item \textbf{Institutional} (four types): modern education (Western philosophy); modern laws (from \textbf{particular to universal} standards --- punishment by crime, not criminal); economic change (jajmani system to modern capitalism: banks, share markets, cooperatives --- sources of mobility increased); political change (democracy, legislature, parliament, Westminster model --- political democracy was a British gift).
\item \textbf{Technological}: printing press, railways, telegraph/telegram, machine-led work (mass production/industry).
\item \textbf{Ideological}: humanitarianism (egalitarianism, meritocracy), rationality, liberty, equality, fraternity, secularism --- leading to social-religious reforms (abolition of sati, gender equality, widow remarriage, girl-child education).
\end{itemize}

Also: a higher standard of life (lifestyle) and \textbf{nationalist consciousness} are effects of westernization.

\begin{QuickRevision}
\begin{itemize}
\item Westernization = change at 3 levels (IIT): Institutional, Ideological, Technological.
\item Institutional: education, laws (particular to universal), economy, politics.
\item Technological: press, railways, telegraph, industry. Ideological: humanitarianism, rationality, etc.
\end{itemize}
\end{QuickRevision}

\section{Negative Effects of Westernization}

Srinivas is \textbf{ambivalent} about westernization (like Weber on bureaucracy) --- he notes both positive and negative effects. Negative effects:
\begin{itemize}
\item \textbf{Rise of individualism} --- weakening of respect for elders, sacrifice, arranged marriage (marriage age now ~26 for men, ~23 for women in urban India); job-switching, live-in relationships.
\item \textbf{Disintegration of the family} --- joint to nuclear (urbanization, migration to new pockets).
\item \textbf{Junk-food culture} --- dietary change, obesity.
\item \textbf{Rise of protest and backward-class movement} --- egalitarianism leading to positive discrimination leading to reservation leading to new communities demanding reservation (Mandal Commission), disintegrating rather than integrating India.
\end{itemize}

\begin{CriticalPoint}
Why does a sociologist call these ``negative''? Because Srinivas is a \textbf{structural functionalist}, concerned with social order and the maintenance of institutions.
\end{CriticalPoint}

\begin{QuickRevision}
\begin{itemize}
\item Negative effects: individualism, family disintegration, junk food, backward-class movement.
\item Srinivas is ambivalent (like Weber on bureaucracy); a structural functionalist concerned with order.
\end{itemize}
\end{QuickRevision}

\section{Primary, Secondary \& Tertiary Westernization}

The impact of westernization was \textbf{not homogeneous} --- it was limited to a few elites (the \textbf{new Indian elites}). Srinivas gives prefixes to westernization:
\begin{itemize}
\item \textbf{Primary westernization} --- those in \textbf{direct contact} with the British: (1) traders/industrialists (Bombay = Parsis, Kolkata = Banias, Madras = Brahmins) --- externally westernized (coat, tie, cigar), internally Sanskritized; (2) intellectuals (Raja Ram Mohan Roy, I. C. Vidyasagar) --- externally Sanskritized, internally westernized (rationality, equality, liberty).
\item \textbf{Secondary westernization} --- those dependent on the primary elites (employees in shipping, patients in hospitals, litigants in courts, readers of newspapers/books) --- consumers, not producers, of knowledge/jobs.
\item \textbf{Tertiary westernization} --- the least westernized: rural folk culture, especially lower castes with remote and limited exposure.
\end{itemize}

\begin{KeyConcept}
Brahmins and Baniyas were the \textbf{``neck points of the funnel''} of westernization --- Brahmins facilitated knowledge, Baniyas facilitated economic growth. They were already privileged (apex of the knowledge pyramid; already in trade) --- no merit, just traditional capital exchanged.
\end{KeyConcept}

\begin{QuickRevision}
\begin{itemize}
\item Impact not homogeneous; limited to new Indian elites.
\item Primary (direct contact), secondary (dependent on primary), tertiary (rural folk/lower castes).
\item Brahmins \& Baniyas = neck points of the funnel.
\end{itemize}
\end{QuickRevision}

\section{Westernization vs Modernization}

\begin{ComparisonBox}
\begin{itemize}
\item \textbf{Daniel Lerner} prefers ``modernization'': ``westernization'' is too local and limited (Japan and Israel are not West; importing from them is modernization, not westernization). Modernization is \textbf{global}.
\item \textbf{M. N. Srinivas} prefers ``westernization'': ``modernization'' is \textbf{ethically loaded} (modern = good), while ``westernization'' is ethically neutral (just drawing from Britain/Europe, no good/bad judgement).
\end{itemize}
\end{ComparisonBox}

\begin{QuickRevision}
\begin{itemize}
\item Lerner: westernization too local/limited; modernization is global.
\item Srinivas: modernization ethically loaded; westernization ethically neutral.
\end{itemize}
\end{QuickRevision}

\section{Westernization \& Sanskritization: A Dynamic Relation}

Sanskritization = lower castes imitating the upper castes; westernization = upper castes imitating the Western/British. Srinivas analysed the dynamic relation between these two. There are \textbf{three relations}:
\begin{itemize}
\item \textbf{Sanskritization as a prerequisite for westernization} --- the primary beneficiaries of westernization were the upper castes (Brahmins, Baniyas), so Sanskritization helps westernize individuals; the two pattern variables are \textbf{imbricated} (conjoined) --- in India there is no binary.
\item \textbf{Westernization promotes/accelerates Sanskritization} --- TV and cinema popularised Sanskritic gods (Krishna, Ramayana, Mahabharat); Harikatha spread via radio/TV; railways and the PRASAD scheme made pilgrimage accessible; the press made epics cheap and accessible; Brahma Samaj, Arya Samaj and Prarthana Samaj popularised the Vedas.
\item \textbf{Conflict between the two} --- primary westernization of the upper castes broadened the hierarchy and triggered backward-class assertion (the \textbf{self-respect movement} of Periyar/E. V. Ramasamy Naicker in South India, leading to anti-Brahminical Dravidian politics, linguistic and cultural consciousness). To be westernized is to no longer respect Brahmins as authority.
\end{itemize}

\begin{CriticalPoint}
Srinivas's conclusion: the relation between Sanskritization and Westernization is \textbf{highly dynamic} and should be \textbf{empirically studied} (not generalised) --- because India is diverse, and each pocket has a different experience of the same seemingly homogeneous concept.
\end{CriticalPoint}

\begin{ExampleBox}
During lockdown, the first thing on television was the Ramayana and Mahabharata --- western technology (TV, the internet) carrying Sanskritic content to the masses.
\end{ExampleBox}

\begin{QuickRevision}
\begin{itemize}
\item Three relations: Sanskritization as prerequisite; westernization promotes Sanskritization; conflict between the two.
\item Conclusion: dynamic relation, empirically studied, not generalised.
\end{itemize}
\end{QuickRevision}
